%=====================================================================
\part{Estrutura Eletrônica de Pontos Quânticos}
%=====================================================================

\noindent
Encerrada a preparação, começa a física. Na Parte~0 tratamos átomos isolados ---
sistemas com um punhado de elétrons, resolvidos em segundos. Agora daremos o
salto que define este livro: sair do átomo e construir um \textbf{ponto quântico},
um pedaço de matéria com dezenas ou centenas de átomos, pequeno o bastante para
que as leis da mecânica quântica moldem suas propriedades de forma dramática.

\vspace{0.3cm}
\noindent
O Capítulo~3 constrói a \emph{geometria}: aprenderemos por que o tamanho de um
nanocristal altera seu \emph{gap}, como recortar uma esfera da rede cristalina do
silício, como curar suas ligações pendentes com hidrogênio e como inserir o átomo
de fósforo que dará origem ao nosso \qubit{}. O Capítulo~4 submeterá essa
geometria a um cálculo Hartree-Fock e extrairá o orbital do doador; o Capítulo~5
refinará o quadro com a Teoria do Funcional da Densidade.
